\chapter{Evaluation}\label{ch:eval}

To evaluate the performance of our guidance algorithm we need to test the two requirements we have set for it: extremeness and realism.

The first concern is whether the algorithm is able to achieve the target, and how well it does so.
This entails testing the algorithm's ability to drive the gap to zero, and how well it does so across different settings.
We will test the algorithm for convergence, stability, and trajectory deviation.
These tests adress will help to unveil potential failure modes of the algorithm and showcase which algorithm design choices are more adequate for different settings.

The second concern is whether the generated states are realistic, in terms of temporal, spatial, vertical, and cross-variable coherence.


--- redundant second sentence, to merge
In our experiments we will test these properties analyzing the impact of the single parameters of our experimental setting.
In particular we will test different modes or ranges for the single parameters and then assess these properties and report on them.

--- ass covering part if I don't manage to perform the realism quantitative tests
In this work we focus on designing a simple guidance method and empirically test its parameters and the properties of the results generated across different settings.
We put less emphasis on the physical realism of the definition of the guidance target and the realism of the generated trajectories,
and defer this to a separate line of work.
Though, we have these two in mind when designing the guidance method throughout.

In addition to these tests, in the discussion we will also discuss different aspects of the results, such as the diversity of the generated trajectories, and the informativeness of the gradient, and other aspects that are not directly tested in the evaluation framework, but are of interest to the overall evaluation of the method and its potential applications.

\section{Flow-time assessment}

\subsection{Parameters}

In flow-time tests we assess:
\begin{itemize}
    \item The algorithm variants
    \item Different guidance schedules
\end{itemize}
with the purpose of selecting the most promising method to perform the final experiments with.

The selection is based on the following criteria.

\subsection{Tests}

Before testiong whether our algorithm produces realistic results we have to check that it actually works.
That is, it achieves the set target, it does so reliably across different settings, and it does so in a stable manner.

In particular we are interested in testing following:

\begin{itemize}
    \item The algorithm converges reliably.
    \item The algorithm doesn't have an oscillatory behaviour in terms of the average trajectory, and in terms of the flow angles and magnitudes of the vector field.
    (the flow trajectory doesn't deviate too much from the unguided trajectory).
\end{itemize}

--- merge the two paragraphs in a single enumeration, can then refer the test by numbers, give the test a clear identifier

--- remove the across runs from the bullets
We will tests these properties with following metrics across different runs:
\begin{itemize}
    \item The gap to the target is driven to zero, and we report the average and std of the gap across different runs.
    \item The flow angles and magnitudes are reported across the flow steps, and we report the average and std of the angles and magnitudes across different runs.
    \item The flow trajectory is compared to the unguided trajectory in PCA latent space, and we report the average and std of the distance between the two trajectories across different runs.
\end{itemize}

Method selection criteria:
\begin{itemize}
    \item Total amount of guidance applied across the flow steps
    \item The pushback from the unguided trajectory
    \item Computational and memory requirements
    \item The number of hyperparameters to tune, and the difficulty of tuning them
    \item The number of runs required to achieve the target across different settings 
    \item The guidance method is robust to different settings, in the sense that it can achieve the target across different settings, and it does so in a stable manner.
\end{itemize}

It's here that we test the developed algorithms and their different variants across different parameter settings, and select the most promising one for the final experiments.


\section{Weather time tests}

\subsection{Parameters}

In an initial phase of the evaluation we want to test parameters in isolation, to both describe their properties, but also to select relevant settings for larger experiments.
We can test these changes by combining them in bigger experiments, that enable us to use other parameters as control.

We start by testing different parameters of the region of interest such as mask centers and sizes across different start states.
We then examine how reproducible are the generated target variable patterns by pushing its most correlated variable instead.
We finally proceed to test intervention across different intensity profiles. 

Throughout these initial experiments we will change start states while setting N=2 and M=3, such that we can control for effects produced by setting N=1,
and effects related to the specific seed, while keeping computational and memory requirements lowest possible.

These experiments will help us understand the properties of the guidance algorithm we developed, and how its parameters influence the results.
As for the guidance method, the described quantitative tests will help us in selecting parameters that produce the most desirable results under this lense.

\subsubsection{Region of interest}
\begin{itemize}
    \item different sizes, same center
    \item different center, same size
    \item different regions across globe
\end{itemize}


\subsubsection{Variable of interest}
For reasons of simplicitly we explicitly only on heatwaves here. 
However, to assess the gradient informativeness we push also other variables in the attempt to reproduce the same pattern on the first experiment.
This might also open up a way to produce more diverse results. 
\begin{itemize}
    \item Temperature first and then push corresponding top-k variables to previosly generated values.
\end{itemize}

\subsubsection{Intervention profile}
\begin{itemize}
    \item Push to different intensities
    \item Assess what happens when you stop pushing (do states reconverge to unguided trajectory?)
\end{itemize}


\subsection{Tests}

Once we have selected the guidance method and its parameters, we want to assess the realism, in addition to other properties, of the generated trajectories.
In particular we want to test the following properties.

\subsubsection{Method tests}
\begin{itemize}
    \item The gradient is informative, in the sense that if we push a not secondary target variable to a certain distribution / average value, we would also expect that if we push this variable to the achieved average value, that the previously targeted variable would achieve the same result.
    \item Diversity test: the generated trajectories are diverse, in the sense that they are not all the same, and they cover a range of possible outcomes.
\end{itemize}

\subsubsection{Realism tests}

Realism tests are of high importance in this work, since we are explicitly guiding the generative model possibly outside of its training distribution.
However, these tests are not particular to this work.
In fact they concern the general problem of evaluating the coherence of data driven models.
That is, the generated states' properties should be consistent with the ground truth dataset, both in the guided and unguided settings.

We design set of tests to assess the coherence to the ground truth dataset of the output trajectories across following dimensions:

\begin{itemize}
    \item temporally (transitions across the same variable should be smooth)
    \item spatially (region of interest)
    \item vertically (pressure levels)
    \item across variables (patterns should be coherent to the learned weather transition function)
\end{itemize}

We will benchmark these properties against the ground-truth data (anchor) and measure the difference between unguided, guided-unguided, and guided states.
For the first three concerns we compute difference maps in the axis of interest and measure min, mean, max, and std across maps.
For the last property we measure all the cross-variable distributions over the weather steps of the ground truth data, and then report the deviation from the average of the unguided, guided-unguided, and guided states.

For all four tests we average over different experiments and report the avg and std over all runs.
This will hopefully give us an indication of whether the generated extremes are faithful to the data distribution.

(Remember, in order for this to work the base-forecast have to have missed potential heat spots --> the extreme is not only about the average, wich the forecast will most probably cover faithfully, but about the extreme spots (max). Both of to be tested anyway.)

Here a visualization summarizing these tests:
--- commented out for speed of compilation
% \begin{figure}[H]
%     \centering
%     \includegraphics[width=0.2\textwidth]{figures/weather-time-tests.png}
%     \caption{Weather-time tests.}
%     \label{fig:weather-time-tests}
% \end{figure}


In the next chapter we will expose the settings for the final experiments, and present the relevant results, and then proceed to discuss them in the following chapter.